\paragraph{Transition-defined policy.} Here each correlation uses the per-boundary
completion levels that make direct-transition compliance rise from 90\% to 95\%:
$.683772\to.776393$ for $\rho=0$ and $.829180\to.908392$ for $\rho=.5$ (the code
uses the unrounded roots). Prediction 3 held: at $e=1$, the full-ICU limits
were 52.36\% and 52.48\%, a contrast of 0.121 percentage points rather than
$-21.5$ under the per-boundary policy. It was a low-risk prediction: its
5-point tolerance is about 40 times the observed contrast, and at $e=1$ the
transition miss probability equals $q_2$ (Section~\ref{sec:s-hygiene}), so
matching transition compliance at both levels matches two-opportunity survival,
the observable that carries $\rho$ in Proposition~\ref{prop:q12}, and the
isolated contrast is exactly zero. The new information is the small ICU
departure from zero and, since transition compliance does not depend on $e$,
the $e=.8$ contrast below. The exactly zero contrast of the isolated
two-opportunity formula is therefore a useful approximation here, not an
identity for the ICU coefficient. At $e=.8$, the full-ICU contrast remained
7.446 points (34.18\% against 41.62\%). Matching the model's transition-compliance target largely removes
the perfect-clearing contrast but does not generally remove dependence on
clustering when clearing is imperfect. The transition-defined policy is not the
same intervention as the per-boundary .90$\to$.95 policy: its per-boundary
completion levels differ by profile, so only the comparison across $\rho$
within one definition is like-for-like, and the forecast depends on which
compliance definition sets the policy target. Neither definition is an
artifact; boundary-scored electronic monitoring records per-boundary events,
and WHO direct observation scores a direct patient-to-patient move as one
opportunity (Section~\ref{sec:s-hygiene} gives the direct-move condition). The construction assumes the stationary
Markov law and one exit and one entry draw per direct move. This is a policy
comparison within the stationary Markov model, not a validation of clinical
audit practice.

